%\documentclass[12pt,oneside]{amsart}
\documentclass{scrartcl}

%\documentclass{article}

\usepackage{amsthm,amsmath,amssymb}
\usepackage[numbers]{natbib}

%\usepackage[colorlinks]{hyperref}
\usepackage{hyperref}
\hypersetup{
    colorlinks,%
    citecolor=black,%
    filecolor=black,%
    linkcolor=black,%
    urlcolor=black
}
\usepackage{hypernat}
\usepackage[top=2.5cm, bottom=2.5cm, left=2.8cm,
right=2.8cm]{geometry}
\usepackage{graphicx}


%\setlength{\parskip}{0pt}
%\setlength{\parsep}{0pt}
%\setlength{\headsep}{0pt}
%\setlength{\topskip}{0pt}
%\setlength{\topmargin}{0pt}
%\setlength{\topsep}{0pt}
%\setlength{\partopsep}{0pt}

\linespread{1}

%\usepackage{marvosym}

%\textwidth = 6.5 in \textheight = 9.2 in \oddsidemargin = 0.0 in
%\evensidemargin = 0.0 in \topmargin = -0.0 in \headheight = 0.0 in
%\headsep = 0.0in \parskip = 0.0in \parindent = 0.0in
%\def\baselinestretch{0.871}


\newtheorem{theorem}{Theorem}[section]
\newtheorem{proposition}[theorem]{Proposition}
\newtheorem{problem}[theorem]{Problem}
\newtheorem{fact}[theorem]{Fact}
\newtheorem{lemma}[theorem]{Lemma}
\newtheorem{defn}[theorem]{Definition}
\newtheorem{corollary}[theorem]{Corollary}
\newtheorem{remark}{Remark}[section]
\newtheorem{example}[theorem]{Example}

\newcommand{\E}{{\mathbb E}}
\newcommand{\se}{{\mathcal{E}}}
\newcommand{\R}{{\mathbb R}}
\renewcommand{\P}{{\mathbb P}}
\newcommand{\ac}{{\mathcal{A}}}
\newcommand{\A}{{\mathcal{A}}}
\newcommand{\B}{{\mathcal{B}}}
\newcommand{\C}{{\mathcal{C}}}
\newcommand{\M}{{\mathcal{M}}}
\newcommand{\Mf}{{\mathfrak{M}}}
\newcommand{\T}{{\mathcal{T}}}
\newcommand{\Z}{{\mathcal{Z}}}
\newcommand{\K}{{\mathcal{K}}}
\newcommand{\lc}{{\mathcal{L}}}
\newcommand{\Ps}{{\mathcal{P}}}
\newcommand{\G}{{\mathcal{G}}}
\newcommand{\pa}{{\mathring{p}}}
\newcommand{\F}{{\cal F}}
\newcommand{\samp}{{\mathcal{X}}}
\newcommand{\X}{{\mathcal{X}}}
\newcommand{\hi}{{\hat{\Phi}}}
\newcommand{\data}{{\text{data}}}

\newcounter{rcnt}[section]
\renewcommand{\thercnt}{(\roman{rcnt})}

\def\qt#1{\qquad\text{#1}}

\def\argmin{\mathop{\rm argmin}}
\def\argmax{\mathop{\rm argmax}}


\setlength{\parskip}{1.4 \medskipamount}
%\sloppy
%\linespread{1.3}

\begin{document}

\title{STAT 238 - Bayesian Statistics  \\
  Lecture Eight} 
\subtitle{Spring 2026, UC Berkeley}

\author{Aditya Guntuboyina}


\date{06 Feb 2026}

\maketitle

The main goal today is to go over Bayesian inference of a parameter $\theta
\in [0, 1]$ from observation $X \sim Bin(n, \theta)$ ($n$ is also
observed) with a Beta distribution prior on $\theta$.

Let us start by recalling the Beta distribution. 

\section{Beta Distribution}
The Beta distribution with parameters $a$ and $b$ is given by the
density: 
\begin{align*}
  f_{\theta}(\theta) = \frac{\theta^{a-1} (1 - \theta)^{b-1}}{B(a, b)} I\{0
  \leq \theta \leq 1\}
\end{align*}
where the normalizing constant $B(a, b)$ is the Beta function given by
\begin{align*}
  B(a, b) = \int_0^1 \theta^{a-1} (1 - \theta)^{b-1} d\theta. 
\end{align*}
For the Beta density to be proper, we need both $a$ and $b$ to be
strictly positive. Here are some basic properties of the Beta
distribution:
\begin{enumerate}
\item When $a = 1$ and $b = 1$, we get the uniform distribution on
  $(0, 1)$. We consider the uniform distribution as the simplest
  example of a Beta distribution.
\item The mean of the Beta distribution is
  \begin{equation*}
    \text{mean} = \frac{a}{a + b}. 
  \end{equation*}
  For example, if $a$ is much smaller compared to $b$, then
  the mean will be small. 
\item The variance of the Beta distribution is
  \begin{equation*}
    \text{variance} = \frac{a}{a + b}
    \frac{b}{a + b} \frac{1}{a + b + 1}. 
  \end{equation*}
 An implication is that, when $a + b$ is large, the variance
 tends to be small, so the Beta density will look skinny.
\item More generally, any moment $\E (\theta^k)$  (for positive integers $k$)
  can be written in closed form as (here $\theta \sim Beta(a, b)$)
  \begin{equation*}
    \E (\theta^k) = \frac{a(a + 1) \dots (a +
      k-1)}{(a + b) (a + b + 1) \dots (a +
      b + k - 1)} \qt{for every $k \geq 1$}. 
  \end{equation*}
\end{enumerate}
In Bayesian inference, improper priors are also sometimes used. In the
case of Beta priors, it is common to use the distributions $Beta(0,
0)$ as a prior. This corresponds to the density
\begin{align*}
  f_{Beta(0, 0)}(\theta) \propto \frac{I\{0 < \theta <
    1\}}{\theta (1 - \theta)}. 
\end{align*}
This improper density cannot be normalized so we don't put any
normalizing constants on the right hand side.

If one wants to interpret $Beta(0, 0)$ as a probability distribution,
the correct answer is the discrete distribution taking
the values 0 and 1 with equal probability:
\begin{align*}
  Beta(0, 0) = Bernoulli(0.5). 
\end{align*}
One way to formalize this equivalence is to argue that the
distribution $Beta(\epsilon, \epsilon)$ converges to $Bernoulli(0.5)$
as $\epsilon \downarrow 0$ in the usual sense of weak convergence. The
easiest way to see this is that the moments of $Beta(\epsilon,
\epsilon)$ converge to the moments of $Bernoulli(0.5)$.

\section{Beta-Binomial Inference}

The goal is to infer $\theta$ from observations $X$ (and $n$) where $X
\sim Bin(n, \theta)$. One can also formulate this problem as that of
estimating $\theta$ from $n$ i.i.d Bernoulli observations $Z_1, \dots,
Z_n \overset{\text{i.i.d}}{\sim} Ber(\theta)$. The likelihood is given
by:
\begin{align*}
  f_{X \mid \theta}(x) = {n \choose x} \theta^x (1 - \theta)^{n-x}
  \propto \theta^{x} (1 - \theta)^{n-x}. 
\end{align*}
The frequentist estimator (MLE) for $\theta$ is $x/n$. For Bayesian
inference, we will use the $Beta(a, b)$ prior:
\begin{align*}
  f_{\theta}(\theta) \propto \theta^{a-1} (1 - \theta)^{b-1} I\{0 \leq
  \theta \leq 1\}
\end{align*}
leading to the posterior:
\begin{align*}
  f_{\theta \mid x}(\theta) &\propto \theta^{a-1} (1 - \theta)^{b-1} I\{0 \leq
                              \theta \leq 1\} \theta^{x} (1 - \theta)^{n-x} \\
  &= \theta^{x + a - 1} (1 - \theta)^{n-x + b - 1} I\{0 \leq \theta
    \leq 1\}. 
\end{align*}
Thus
\begin{align*}
  \theta \mid x \sim Beta(x + a, n-x + b). 
\end{align*}
A natural  point estimate of $\theta$ is the posterior mean:
\begin{align*}
  \hat{\theta} = \E(\theta \mid x) = \frac{x+a}{n + a + b}. 
\end{align*}
Here are two things to note about the posterior mean:
\begin{enumerate}
\item It is a convex combination of the MLE and the prior mean
  $a/(a+b)$:
  \begin{align*}
    \frac{x+a}{n + a + b} = \left(\frac{n}{a+b+n} \right) \frac{x}{n}
    + \left(\frac{a+b}{a+b+n} \right) \frac{a}{a+b}. 
  \end{align*}
  So if $n$ is much larger compared to $a+b$, then the posterior mean
  will be very close to the MLE. Conversely if $n$ is much smaller
  compared to $a + b$, then the posterior mean will be very close to
  the prior mean.
\item When $a = b = 0$, then the posterior mean exactly coincides with
  the MLE. Thus if we use the $Beta(0, 0)$ prior, then posterior
  inference will be close to frequentist inference. Note again that
  this prior is improper.
\item If we use the $Beta(0, 0)$ prior and if we observe $n = x = 1$
  (i.e., we only have data on one Bernoulli trial which resulted in a
  heads). Then the posterior is $Beta(x + a, n-x + b) =
  Beta(1,0)$. $Beta(1, 0)$ is also improper but it can be
  interpreted as the point mass at 1:
  \begin{align*}
    Beta(1, 0) = Bernoulli(1) = \delta_{\{1\}}. 
  \end{align*}
  This can be proved, for example, by computing the moments of
  $Beta(1, \epsilon)$ and showing that they all approach 1 (which are
  the moments of $\delta_{\{1\}}$) as $\epsilon \rightarrow 0$. 
  Similarly for $x = 0$ and $n = 1$, the posterior becomes $Beta(0,
  1)$ which should be interpreted as $\delta_{\{0\}}$. If $n = 2$ and
  $x = 1$ (i.e., we observe one head and one tail), the posterior is
  the uniform distribution on $(0, 1)$.
\item The process of going from the $Beta(0, 0)$ prior to the $Beta(1,
  0)$ posterior when $n = x = 1$ is described as intuitive in the
  following situation by \citet[Section 12.4.3]{Jaynes}: \textit{...in a
    chemical laboratory we find a jar containing an unknown and
    unlabeled compound. We are at first completely ignorant as to
    whether a small sample of this compound will dissolve in water or
    not. But, having observed that one small sample does dissolve, we
    infer immediately that all samples of this compound are water
    soluble, and although the conclusion does not carry quite the
    force of deductive proof, we feel strongly that the inference was
    justified}. 
\end{enumerate}

\section{Multiple Replications of the Same Problem}
Consider the problem of estimating $\theta_i$ from
$(X_i, n_i)$ for each $i = 1, \dots, N$, with $X_i \sim Bin(n_i,
\theta_i)$. For a concrete example, take all counties in the United
States with $X_i$ denoting the number of deaths
due to kidney cancer (in the period $1980-89$) and $n_i$ denoting the
average population (during the same period) for the $i$-th county.

It is easy to see that the naive frequentist estimate $X_i/n_i$ can be
very bad for $\theta_i$ if $n_i$ is small. We will therefore use
Bayes estimation using the prior:
\begin{align*}
  \theta_i \overset{\text{i.i.d}}{\sim} Beta(a, b). 
\end{align*}
The choice of $a$ and $b$ is crucial here. Since we have a large
dataset, it makes sense to learn good values of $a$ and $b$ from the
observed data. We will look at the following three ways of doing
this.

\begin{enumerate}
\item \textbf{Method One}: We place a threshold on $n_i$ and filter out
  $i$ for which $n_i$ exceeds the threshold. We then select $a$ and
  $b$ by fitting the $Beta(a, b)$ density to the data $\{X_i/n_i : n_i
  \geq \text{threshold}\}$. The Beta density fitting can be done by
  just matching mean and variances. Let $m$ and $V$ denote the mean
  and variance of $\{X_i/n_i : n_i
  \geq \text{threshold}\}$. Then we obtain $a$ and $b$ by solving:
  \begin{equation*}
    \frac{a}{a+b} = m ~~~ \text{ and } ~~~ \frac{a}{a+b} \frac{b}{a+b}
    \frac{1}{a+b+1} = V
  \end{equation*}
  which gives
  \begin{equation*}
    \hat{a} = m \left(\frac{m(1-m)}{V} - 1 \right) \text{ and } \hat{b} = (1
    - m) \left(\frac{m(1-m)}{V} - 1 \right). 
  \end{equation*}
  One drawback of this method is that the selection of threshold can
  be arbitrary. In the kidney cancer example, we choose 300000 as the
  threshold, but we could have also chosen some other threshold such
  as 350K or 400K.
\item \textbf{Method Two}: We consider the marginal likelihood of
  $X_i$ given $a, b$. Note that
  \begin{align*}
    \P \{X_i = x \mid a, b\} &= \int_0^1  \P\{X_i = x \mid \theta_i = \theta\}
                               f_{\theta_i \mid a, b}(\theta) d\theta
    \\
    &= \int_0^1 {n_i \choose x} \theta^x (1 - \theta)^{n_i - x}
      \frac{\theta^{a-1} (1 - \theta)^{b-1}}{B(a, b)} d\theta \\
    &={n_i \choose x} \frac{1}{B(a, b)} \int_0^1 \theta^{x + a - 1} (1 -
      \theta)^{n_i - x + b - 1} d\theta \\ &= {n_i \choose x} \frac{B(x + a,
      n_i - x + b)}{B(a, b)}
  \end{align*}
  As a result, the marginal likelihood of the data  given  in terms of
  $a$, $b$ is:
  \begin{align*}
    \prod_{i=1}^N {n_i \choose X_i} \frac{B(X_i + a,
      n_i - X_i + b)}{B(a, b)}
  \end{align*}
  So the maximum likelihood estimates of $a$, $b$ are given by:
  \begin{align*}
    (\hat{a}_{\text{MLE}}, \hat{b}_{\text{MLE}}) = \underset{a,
    b}{\text{argmax}} \prod_{i=1}^N {n_i 
    \choose X_i} \frac{B(X_i + a, n_i - X_i + b)}{B(a, b)}.  
  \end{align*}
  This maximization can be done numerically by taking a grid of $a$
  and $b$ values.
\item \textbf{Method Three}: This is the fully Bayes solution where we
  simply place priors on $a$ and $b$. It is sensible to reparametrize
  as:
  \begin{align*}
    \mu = \frac{a}{a+b} ~~ \text{ and } ~~ \kappa = a + b. 
  \end{align*}
  Now we take the priors:
  \begin{align*}
    \mu \sim \text{uniform}(0, 1) ~~ \text{ and } ~~ \log \kappa \sim
    \text{uniform}(-\infty, \infty). 
  \end{align*}
  Inference under this fully Bayesian model is nontrivial from a
  computational point of view. We shall discuss this in the next
  lecture. 
\end{enumerate}




\begin{thebibliography}{99}

  
%\bibitem[Fisher(1934)]{Fisher1934}
%R.~A. Fisher,
%\newblock ``Probability, likelihood and quantity of information in the logic of
%uncertain inference,''
%\newblock \emph{Proceedings of the Royal Society of London. Series~A,
%Containing Papers of a Mathematical and Physical Character}, vol.~146,
%no.~856, pp.~1--8, 1934.

     \bibitem[Jaynes(2003)]{Jaynes} Jaynes, E. T, (2003)
  \textit{Probability theory: the logic of science}. Cambridge
  University Press, 2003.

%  \bibitem[Sinai(1992)]{Sinai} Sinai, Y. G, (1992)
%  \textit{Probability theory: an introductory course}. Springer
%  Verlag, 1992.  

%\bibitem[{deGroot(1986)}]{DeGrootBlackwell}DeGroot, M. H. (1986)
%       A conversation with David Blackwell.
%\newblock \emph{Statistical Science}, \textbf{1}, 40--53.

%         \bibitem[Mosteller(1987)]{Mosteller} Mosteller, F, (1987)
%  \textit{Fifty challenging problems in probability with
%    solutions}. Courier Corporation, 1987.

%    \bibitem[MacKay(2003)]{MacKay} MacKay, D. J. C., (2003)
%  \textit{Information theory, inference, and learning
%  algorithms}. Cambridge University Press, 2003.

%\bibitem[{Lindley(2013)}]{Lindley}Lindley, D. V. (2013)
%   \textit{Understanding Uncertainty}. John Wiley and sons, 2013.


\end{thebibliography}







\end{document}

%%% Local Variables:
%%% mode: latex
%%% TeX-master: t
%%% End:
